% ============================================================
% 3. Transparencias de distribución en RabbitMQ (Sánchez)
% ============================================================
\section{Transparencias de distribución}

\begin{frame}{Transparencias presentes en la arquitectura}
  \small
  \begin{tabularx}{\textwidth}{@{}p{2.5cm}p{3.5cm}X@{}}
    \toprule
    \textbf{Transparencia} & \textbf{Qué oculta} & \textbf{Ejemplo en el diagrama} \\
    \midrule
    Acceso & \raggedright Diferencias en cómo se representan y acceden los
      recursos. &
      La tienda y los trabajadores solo comparten el protocolo AMQP, no su
      implementación \cite{rabbitmq2026connections}. \\
    Ubicación & \raggedright Dónde está el recurso. &
      La tienda publica en el \textit{exchange} \texttt{pedidos} por su
      nombre; no sabe qué trabajador ni qué nodo atiende el pedido
      \cite{rabbitmq2026amqp,rabbitmq2026clustering}. \\
    Replicación & \raggedright Que existen varias copias. &
      \texttt{procesar\_pedidos} tiene réplicas en los 3 nodos, pero se usa
      como una sola cola \cite{rabbitmq2026quorum}. \\
    Concurrencia & \raggedright Que otros comparten el recurso. &
      Los tres trabajadores consumen de la misma cola sin coordinarse entre
      sí \cite{rabbitmq2026amqp}. \\
    Fallos & \raggedright El fallo y la recuperación. &
      Si el trabajador 3 cae sin \texttt{ack}, su pedido pasa a otro; si cae
      el nodo líder de la cola, se elige otro
      \cite{rabbitmq2026queues,rabbitmq2026quorum}. \\
    \bottomrule
  \end{tabularx}

  \vspace{0.2cm}
  {\tiny Tipos de transparencia: \cite{vansteen2023distributed}.}
\end{frame}

\begin{frame}{Limitaciones de esas transparencias}
  \begin{itemize}
    \item \textbf{Fallos:} el pedido reentregado llega marcado como
          \texttt{redelivered} y puede cambiar el orden; si el trabajador
          alcanzó a procesarlo, se procesaría dos veces
          \cite{rabbitmq2026queues,fu2021fair}.
    \item \textbf{Ubicación:} si cae el nodo al que está conectada la
          tienda, esta debe reconectarse a otro nodo
          \cite{rabbitmq2026clustering}.
    \item \textbf{Replicación:} por defecto, una cola vive en un solo nodo;
          la de quórum necesita mayoría, así que con 3 nodos tolera 1 fallo
          \cite{rabbitmq2026clustering,rabbitmq2026quorum}.
    \item \textbf{Latencia:} la tienda no sabe cuándo se procesará el
          pedido, y la demora de la red no se puede ocultar
          \cite{vansteen2023distributed}.
  \end{itemize}

  \begin{block}{Conclusión}
    La transparencia total no es alcanzable y tiene un costo en rendimiento
    \cite{vansteen2023distributed}. RabbitMQ oculta la distribución en el
    caso normal, pero ante fallos y latencia la aplicación debe intervenir.
  \end{block}
\end{frame}
